\section{INTRODUCTION}

Chinese equipment and military-technology documents contain equipment names,
modules, technical schemes, and performance metrics that are valuable for
downstream search, analysis, and knowledge-base construction. Manual annotation
of these entities is expensive, particularly when a domain corpus grows beyond
a small collection of passages. Few-shot NER provides a practical route to
bootstrap extraction from limited labeled examples, but generative LLM
extractors can also create entity-level false positives. For example, an
organization may be labeled as \entitytype{EQUIPMENT}, a descriptive phrase as
\entitytype{MODULE}, or a partial numerical expression as a complete
\entitytype{METRIC}. Such errors increase the burden of downstream review and
reduce trust in a deployment.

The practical question in this setting is therefore not only how to maximize
F1, but also how to reduce unsupported or boundary-unreliable entity
predictions without materially degrading extraction performance. Few-shot NER
studies have shown that example selection and semantic proximity affect
performance under limited supervision \cite{ding2021fewnerd,yang2020structured,ma2022decomposed}.
Meanwhile, universal information extraction and instruction-tuned extraction
systems show that structured generation can support multiple extraction tasks
\cite{lu2022uie,wang2023instructuie,wang2023gptner,luo2025oneke}. These lines of
work motivate a targeted engineering question: can relevant demonstrations,
query-matched terminology, and retrieved evidence provide useful and
interpretable control over a local LLM extractor in a narrow equipment-domain
NER setting?

This paper studies that question using a frozen Chinese equipment-domain
dataset and four incremental configurations. Configuration A uses a
fixed-random demonstration set; B replaces it with kNN-retrieved
demonstrations; C adds a query-matched domain glossary; and D additionally
injects reranked RAG evidence. The resulting comparison chain is deliberately
narrow. It measures entity-level micro precision, recall, F1, and
false-positive entity counts for NER only. It does not establish
relation-extraction effectiveness, Chain-of-Verification (CoVe) gains,
public-benchmark generalization, $K$ sensitivity, statistical significance, or
completion of the originally planned full information-extraction matrix.

The contributions are as follows.

\begin{enumerate}
  \item We document a small frozen Chinese equipment-domain NER dataset with
  four entity types and an independent two-human non-test calibration audit.
  The audit is reported accurately as non-test calibration agreement, not as
  formal frozen-test inter-annotator agreement (IAA).
  \item We implement a reproducible local-LLM few-shot extraction pipeline
  that compares fixed-random versus kNN demonstrations, query-matched glossary
  injection, and reranked RAG evidence.
  \item We provide descriptive three-seed evidence that kNN and the glossary
  improve F1/recall, while RAG preserves near-tied F1 relative to the glossary
  configuration and improves precision with 55.0 fewer false-positive entities
  on average. No claim of statistical significance is made.
\end{enumerate}
